\section{Conclusion}
\label{sec:conclusion}

We evaluated whether post-training concept erasure reproduces the behavior of
matched models trained without loss on concept-associated chunks. The answer
depends on which property is measured. RMU was closest to the twins in signed
mean target-answer NLL for all three concepts, but the question-level analysis
showed that cancellation explained the apparent Baseball and AI matches. SNMF
had the lowest target KL and was closer than the full model on 64--66\% of
target questions, while still producing little target erasure. EMBER produced
the strongest target-specific NLL increase and the highest combined
efficacy--preservation score for Baseball and AI; RMU obtained the highest
score for Rome. These differences also show that method rankings vary across
concepts and evaluation criteria.

Most notably, successful target suppression with preserved neighboring and
SciQ accuracy did not imply similarity to the concept-excluded twin. EMBER
had the highest $H_{\mathrm{test}}$ for Baseball and AI while remaining the
most distant erasure method from their twins under KL and both signed-mean and
question-level NLL comparisons. SNMF showed the complementary disagreement:
it was consistently closest to the twin but produced little target erasure.
Efficacy--preservation and similarity to a
concept-excluded twin therefore capture distinct properties and should be
reported separately. This conclusion is limited to the tested models, three
concepts, fixed-answer evaluations, and one training seed. A useful next step
is to repeat the comparison across additional concepts and seeds to determine
how consistently the observed disagreement holds.
